\documentclass[11pt]{article}
\usepackage[a4paper,margin=2.4cm]{geometry}
\usepackage{amsmath,amssymb}
\usepackage{booktabs}
\usepackage[hidelinks]{hyperref}
\usepackage{xcolor}
\newcommand{\code}[1]{\texttt{\small #1}}

\title{The Poisson Mode of GBDTE:\\ What It Is, How It Works, and What Is Broken}
\author{Package documentation --- \code{golang/poisson\_legacy} and\\
\code{python/extra\_boost\_py/experiments/poisson\_bench.py}}
\date{July 3, 2026}

\begin{document}
\maketitle

\begin{abstract}
This document explains, from first principles, what the Poisson mode of the GBDTE
package models, how the legacy Go implementation trains it, how the new synthetic
benchmark generator produces data for it, and --- based on the empirical validation of
2026-07-02 --- where the implementation is inconsistent and why. It is meant as the
shared mental model before fixing the engine. Statements about the code cite exact
files and functions; statements about behavior cite the reproducible tests in
\code{python/tests/test\_poisson\_semantics.py}.
\end{abstract}

\section{The problem the Poisson mode solves}

\subsection{Counts instead of values or labels}

The package's other two losses answer ``what value?'' (MSE) and ``which class?''
(LogLoss). The Poisson mode answers \emph{``how many events, and at what rate?''}:
insurance claims per policy-year, machine failures per month, calls per hour, air-traffic
conflicts per sector-day. The classical tool is Poisson regression, a generalized linear
model in which a count $y$ is modeled as
\[
y \sim \mathrm{Poisson}(\mu), \qquad \mu = \mathbb{E}[y \mid x],
\]
and the model is fitted by maximizing the Poisson log-likelihood
$\sum_i \left( y_i \log \mu_i - \mu_i \right)$ (up to a term not depending on the model)
\cite{mccullagh1989,cameron2013}. Modern gradient-boosting packages all offer this as an
objective: LightGBM (\code{objective=poisson}), XGBoost (\code{count:poisson}), CatBoost
(\code{loss\_function=Poisson}) \cite{lgbmdocs,xgbdocs,cbdocs}; in actuarial practice
Poisson-boosted trees are the standard workhorse for claim frequency
\cite{ohlsson2010,denuit2020,wuthrich2023,yang2018}.

\subsection{The GBDTE twist: intensity that drifts in time}

All of the above predict a \emph{constant} expected count per feature vector. GBDTE's
premise is different: the event rate is a \emph{function of time},
\[
\lambda(t \mid x) \;=\; \text{events per unit time for an object with features } x,
\]
and what we care about is \emph{extrapolating} $\lambda$ into a future time window that
the training data never saw. This is the inhomogeneous Poisson process setting
\cite{daley2003}: over a window $W$, an object with intensity $\lambda(t)$ produces
$\mathrm{Poisson}\!\left(\int_W \lambda(t)\,dt\right)$ events, and the log-likelihood of
observed event times $t_1,\dots,t_n$ is
\begin{equation}
\label{eq:pp-loglik}
\ell \;=\; \sum_{i=1}^{n} \log \lambda(t_i) \;-\; \int_W \lambda(t)\,dt .
\end{equation}

GBDTE applies its usual role separation to this problem:
\begin{itemize}
\item \textbf{partition features} $X^{(p)}$ (stable object attributes) drive the tree
      splits and identify \emph{which} intensity regime an object belongs to;
\item \textbf{extrapolation basis} $\varphi(t)$ (e.g.\ $(1, t)$) lives inside each leaf:
      the leaf stores a weight vector $w$ and predicts the time-varying intensity
      \[
      \lambda(t) \;=\; w^\top \varphi(t)
      \]
      --- a \emph{linear intensity model per leaf}, so the fitted rate has a trend that
      can be carried into the future window.
\end{itemize}
A conventional Poisson booster puts $t$ among the ordinary features and therefore
predicts a piecewise-constant function of $t$; beyond the last training split it
extrapolates a constant. GBDTE's leaf-linear intensity is exactly what
distinguishes it --- the same argument as in the MSE and LogLoss modes, applied to
event rates. Note the additive form $\lambda = w^\top\varphi(t)$ differs from the
GLM convention $\log\mu = x^\top\beta$ \cite{mccullagh1989}: the intensity is linear,
not log-linear, which permits a closed-form Newton step per leaf but requires care to
keep $\lambda > 0$.

\section{The data contract of \code{PoissonLegacyBooster}}

The Python API (\code{python/extra\_boost\_py/poisson\_booster.py}) takes five arrays.
The semantics below were validated empirically on 2026-07-02
(\code{test\_poisson\_semantics.py}); the engine itself documents none of this.

\begin{center}\small
\begin{tabular}{lll}
\toprule
Argument & Shape & Meaning \\
\midrule
\code{bjids} & $(n,)$ int32 & object (cohort) identifier; rows of one object share it \\
\code{freqs} & $(n,)$ f64 & event \textbf{count} observed in the row's time bin; must be $> 0$ \\
\code{features\_inter} & $(n, p)$ f64 & partition features (constant within an object) \\
\code{features\_extra} & $(n, d)$ f64 & $\varphi(t_{\text{row}})$ evaluated at the row's bin center \\
\code{psi} & $(d,)$ f64 & $\int_{W} \varphi(t)\,dt$ over the \emph{whole observation window} \\
\bottomrule
\end{tabular}
\end{center}

So one ``row'' is \emph{(object, time bin, count)}: an object observed over $W=[0,c]$
partitioned into bins contributes one row per bin. $\psi$ carries the exposure term of
Eq.~\eqref{eq:pp-loglik}: for a leaf model $\lambda = w^\top\varphi$,
$\int_W \lambda = w^\top \psi$ per object.

Hard constraints found empirically:
\begin{itemize}
\item \textbf{zero counts crash the process} (a \code{munmap\_chunk} abort inside the
      CGO layer) --- the benchmark generator therefore guarantees counts $\ge 1$;
\item \code{predict(F\_inter, $\varphi$(t))} returns the intensity \textbf{in per-bin
      count units}, i.e.\ the expected count of one row, not a rate per unit time;
\item only \code{n\_stages=1} produces consistent extra-mode results (Section~\ref{sec:bug}).
\end{itemize}

\section{How training works (\code{golang/poisson\_legacy})}

Training alternates tree building and bias accumulation
(\code{booster.go:Train}): each stage builds one tree against the current prediction
(\emph{bias}), then adds the tree's predictions to the bias. Leaf values of the first
tree are used at full scale; later trees are damped by the learning rate
(\code{tree.go:leafFromPrediction}). The interesting part is how leaf weights are
computed, and it differs between the first tree and all subsequent ones
(\code{pmatrix.go:WholeLoss} dispatches on \code{bias == nil}).

\subsection{First tree, extra mode: least squares on counts}

\code{wholeLossFirstTreeExtra} solves, for every candidate leaf $L$,
\begin{equation}
\label{eq:first-tree}
\Big( \sum_{i \in L} \varphi_i \varphi_i^\top + \lambda_{\text{reg}} I \Big)\, w
\;=\; \sum_{i \in L} y_i\, \varphi_i ,
\end{equation}
i.e.\ an \emph{ordinary least-squares regression of the counts on the basis}. This is
not the Poisson likelihood --- it is its Gaussian surrogate --- but it has two convenient
properties: it is exact in one shot (no iteration), and if the counts follow a linear
trend it reproduces it exactly. This is why the engine's own Go test
(\code{TestExtraPoissonFirstTree}) sees per-row predictions equal to the training
frequencies, and why our validation recovered $\lambda_0(t) = 2+2t$ vs
$\lambda_1(t) = 4-2t$ perfectly with one stage. The fitted $w$ is in
\textbf{count-per-bin units} because the targets $y_i$ are counts. Note that
Eq.~\eqref{eq:first-tree} does not use $\psi$ at all.

\subsection{Split search}

For each split candidate the engine evaluates the post-fit loss on both sides using
cumulative-sum tables over rows sorted per feature (\code{Gax}), exactly like the main
GBDTE engine; an optional \code{unbalanced\_penalty} discourages extreme splits. Leaf
weight vectors for the winning split come out of the same pass (\code{DeltaUp} /
\code{DeltaDown}).

\subsection{Later trees, extra mode: Newton on the Poisson process likelihood}

\code{wholeLossNextTreeExtra} treats the new tree as an \emph{additive correction}
$u(t) = w^\top\varphi(t)$ on top of the current intensity $b(t)$ (the bias), and takes
one damped Newton step on Eq.~\eqref{eq:pp-loglik} around $u = 0$:
\begin{equation}
\label{eq:newton}
g \;=\; N_L\, \psi \;-\; \sum_{i \in L} \frac{y_i\, \varphi_i}{b_i},
\qquad
H \;=\; \sum_{i \in L} \frac{y_i\, \varphi_i \varphi_i^\top}{b_i^2},
\qquad
w \;=\; -H^{-1} g,
\end{equation}
where $N_L$ counts distinct objects in the leaf (\code{countObjects}). This is the
correct derivative structure for Eq.~\eqref{eq:pp-loglik}: the exposure term
contributes $N_L \psi$ (each object integrates the correction over the whole window),
the event term contributes $\sum y\varphi/b$.

\section{The inconsistency, precisely (fixed 2026-07-04)}
\label{sec:bug}

\noindent\fbox{\parbox{\linewidth}{\small \textbf{Resolved.} The defect described in this
section was fixed on 2026-07-04 by adopting Fix~1 below: the Newton exposure term now uses
the per-row running sum $\sum_i\varphi_i$ instead of $N_L\psi$. Multi-stage boosting now
converges to the true intensity; \code{test\_extra\_mode\_multi\_stage\_is\_consistent}
passes for \code{n\_stages} in $\{2,5\}$. The analysis below is kept for the record.}}

\medskip
Equations \eqref{eq:first-tree} and \eqref{eq:newton} do not speak the same units.

\begin{itemize}
\item The first tree fits $b(t)$ in \textbf{per-bin counts}: its stationary condition
      matches $\sum y_i \varphi_i$, where $y_i$ are counts in bins of width $\Delta$.
\item The Newton stage's stationary condition, $\sum_i y_i \varphi_i / \lambda(t_i) =
      N_L\,\psi$, is satisfied by an intensity in \textbf{events per unit time}: for a
      constant true rate $r$, plugging $\lambda = r$ makes the left side
      $\approx N_L \int_W \varphi$ only because
      $\sum_{\text{bins}} (r\Delta)\varphi(t_b)/r = \Delta \sum_b \varphi(t_b) \approx
      \int_W \varphi$. Plugging the first tree's count-scale intensity
      $\lambda = r\Delta$ instead inflates the left side by $1/\Delta$.
\end{itemize}

So after a perfect first tree, the second stage's gradient is far from zero --- not
because the fit is bad, but because the exposure term evaluates it on a different scale.
The observed behavior follows: with \code{n\_stages=1} the recovery of a known two-group
linear intensity is exact (to sampling noise); with \code{n\_stages=2} every prediction
shifts by a systematic offset; with many stages the iteration settles at a fixed point
that is \emph{not} the maximizer of the engine's own likelihood (verified by solving the
MLE independently). This defect was pinned by
\code{test\_extra\_mode\_multi\_stage\_is\_consistent} (formerly a strict \code{xfail},
now a passing test).

\paragraph{Fix directions.} Fix~1 was implemented.
\begin{enumerate}
\item \textbf{Make everything per-row Poisson regression on binned counts.} Replace the
      per-object exposure $N_L\psi$ with a per-row exposure $\varphi(t_i)\Delta_i$
      (row $i$'s bin integrates $\int_{\text{bin}} \varphi \approx \varphi(t_b)\Delta$).
      Then the likelihood is the standard binned-count Poisson likelihood, units agree
      with the first tree automatically, $\psi$ and \code{bjids} become unnecessary,
      zero counts become legal, and results match what external Poisson boosters
      optimize. This is the simplest and most conventional repair
      \cite{cameron2013,ohlsson2010}.
\item \textbf{Keep the point-process form, fix the first tree.} Fit the first tree by
      the same Newton objective (starting from a constant initial intensity) so all
      stages optimize Eq.~\eqref{eq:pp-loglik} in rate units, and document that
      \code{freqs} rows are event/bin counts while predictions are rates. More faithful
      to the ``true'' inhomogeneous process, but more work and still needs a
      positivity guard for $\lambda$.
\end{enumerate}
Either way, a positivity safeguard for leaf intensities matters: the current linear
$\lambda(t)$ fitted on $t < 0.6$ can cross zero inside the evaluation window, which is
what blew up the Poisson deviance at some depths in the benchmark run.

\section{The benchmark generator}

\code{poisson\_bench.py} builds the Poisson analog of the paper's grouped MSE benchmark
--- ``grouped drifting intensity'':

\begin{enumerate}
\item \textbf{Objects are cohorts.} Each of \code{n\_objects} objects represents a
      cohort of \code{exposure} identical units (e.g.\ a portfolio segment). Cohorts
      make per-bin counts comfortably positive, which the engine requires.
\item \textbf{Latent group $\to$ intensity.} Object $j$ carries binary partition
      features encoding a latent group $g$ out of \code{k\_groups} (plus the continuous
      nuisance feature $f_m$). The group's per-unit intensity is linear in time:
      \[
      \lambda_g(t) \;=\; \beta_{g,0} + \beta_{g,1}\, t,
      \qquad \beta_{g,0} \sim U(1,5),\quad
      \beta_{g,1} = \text{drift} \cdot u \cdot \beta_{g,0},\ u \sim U(-0.8, 0.8),
      \]
      so intensities stay positive on $[0,1]$ by construction.
\item \textbf{Role purity $\rho$.} As in the MSE family, the actual intensity is a
      mixture $\rho\,\lambda_{\text{group}} + (1-\rho)\,\lambda_m$, where $\lambda_m$
      depends smoothly on the nuisance $m$; at $\rho = 1$ the partition features carry
      all structure, at $\rho = 0$ none.
\item \textbf{Binned observation.} Time $[0,1]$ splits into \code{n\_bins} equal bins;
      the row for (object $j$, bin $b$) gets count
      $y \sim \mathrm{Poisson}\!\big(\lambda_j(t_b)\, \Delta \cdot \text{exposure}\big)$,
      clipped to $\ge 1$ (measured distortion at the default exposure: the average
      zero-probability is $6\%$; at exposure 60 it vanishes).
\item \textbf{Temporal split.} Train on $t < 0.6$, evaluate on $t \ge 0.6$: the model
      must extrapolate each cohort's trend into an unseen future window --- the same
      protocol as the MSE ($t<0.5$) and LogLoss ($t<0.4$) benchmarks.
\item \textbf{Dual form.} The same dataset is exposed (a) as events
      (\code{bjids, freqs, features, $\varphi$, $\psi$}) for GBDTE and (b) as a plain
      table (rows = object $\times$ bin, target = count, features = partition
      candidates + $t$) for LightGBM/XGBoost/CatBoost with their native Poisson
      objectives. A \code{lam\_true} column stores the true intensity, enabling an
      oracle upper bound and an intensity-RMSE diagnostic that real data cannot offer.
\end{enumerate}

Evaluation uses the mean Poisson deviance on the future window,
\[
D \;=\; \frac{2}{n}\sum_i \Big[ y_i \log \tfrac{y_i}{\hat\mu_i} - (y_i - \hat\mu_i) \Big],
\]
the standard proper scoring rule for count models \cite{cameron2013,wuthrich2023}, plus
the synthetic-only intensity RMSE against \code{lam\_true}.

\section{Where things stand}

With the engine restricted to a single stage, GBDTE-Poisson currently \emph{loses} to
conventional Poisson boosters on this benchmark (mean deviance $1.02$--$1.03$ vs.\
$0.85$--$0.89$ for CatBoost/LightGBM; Bayes oracle $0.79$--$0.80$; 10 seeds, Friedman
$p = 1.3\cdot 10^{-20}$), and the wrong-roles ablation is worst ($1.45$--$1.50$), as
expected. The loss is structural: one tree must simultaneously discover the partition
and fit stable leaf trends, with no boosting to correct residual errors --- and
mis-fitted leaf trends extrapolate badly (occasionally through zero). Fixing the
multi-stage update per Section~\ref{sec:bug} is therefore the prerequisite for any
paper-grade Poisson claim; the benchmark, baselines, statistics, and reports are already
in place and re-run with one command:
\begin{center}
\code{PYTHONPATH=python .venv/bin/python scripts/run\_article\_experiments.py
--suite baselines\_poisson --seeds 10}
\end{center}

\begin{thebibliography}{99}\small

\bibitem{mccullagh1989} P.~McCullagh and J.~A.~Nelder.
\emph{Generalized Linear Models}, 2nd ed. Chapman \& Hall, 1989.

\bibitem{cameron2013} A.~C.~Cameron and P.~K.~Trivedi.
\emph{Regression Analysis of Count Data}, 2nd ed. Cambridge University Press, 2013.

\bibitem{daley2003} D.~J.~Daley and D.~Vere-Jones.
\emph{An Introduction to the Theory of Point Processes, Vol.~I: Elementary Theory and
Methods}, 2nd ed. Springer, 2003. (Inhomogeneous Poisson processes and the
log-likelihood \eqref{eq:pp-loglik}.)

\bibitem{friedman2001} J.~H.~Friedman.
Greedy function approximation: a gradient boosting machine.
\emph{Annals of Statistics} 29(5):1189--1232, 2001.

\bibitem{ohlsson2010} E.~Ohlsson and B.~Johansson.
\emph{Non-Life Insurance Pricing with Generalized Linear Models}. Springer, 2010.
(Poisson claim-frequency modeling with exposure offsets.)

\bibitem{denuit2020} M.~Denuit, D.~Hainaut, and J.~Trufin.
\emph{Effective Statistical Learning Methods for Actuaries II: Tree-Based Methods and
Extensions}. Springer, 2020. (Poisson deviance as the split criterion for count data.)

\bibitem{wuthrich2023} M.~V.~W\"uthrich and M.~Merz.
\emph{Statistical Foundations of Actuarial Learning and its Applications}.
Springer Actuarial, 2023. Open access:
\url{https://link.springer.com/book/10.1007/978-3-031-12409-9}.

\bibitem{yang2018} Y.~Yang, W.~Qian, and H.~Zou.
Insurance premium prediction via gradient tree-boosted Tweedie compound Poisson models.
\emph{Journal of Business \& Economic Statistics} 36(3):456--470, 2018.
arXiv:\href{https://arxiv.org/abs/1508.06378}{1508.06378}.

\bibitem{lgbmdocs} LightGBM documentation: \code{objective=poisson}.
\url{https://lightgbm.readthedocs.io/en/latest/Parameters.html}.

\bibitem{xgbdocs} XGBoost documentation: \code{count:poisson}.
\url{https://xgboost.readthedocs.io/en/stable/parameter.html}.

\bibitem{cbdocs} CatBoost documentation: regression loss functions (Poisson).
\url{https://catboost.ai/docs/en/concepts/loss-functions-regression}.

\end{thebibliography}
\end{document}
